% Figura: problema 5.1(a). Esfera uniformemente cargada: campo, potencial y densidad
% de energía frente a r (unidades Q/(4 pi eps0 a^2), Q/(4 pi eps0 a) y valor en r = a).
\begin{figure}[t]
\centering
\begin{subfigure}[b]{0.28\textwidth}\centering
\begin{tikzpicture}[scale=1.0]
  \shade[inner color=acentoB!45!fondo, outer color=acentoB!18!fondo] (0,0) circle (1.05);
  \draw[texto, line width=0.5pt] (0,0) circle (1.05);
  \draw[gauss] (0,0) circle (0.55);
  \node[etiq, text=acentoB] at (-45:0.8) {$r$};
  \draw[texto, line width=0.4pt, -{Stealth[length=1.5mm]}] (0,0) -- (40:1.05);
  \node[etiq] at (40:1.32) {$a$};
  \node[etiq] at (0,-1.55) {$\rho=\dfrac{3Q}{4\pi a^{3}}$};
\end{tikzpicture}
\caption{Geometría}
\end{subfigure}\hfill
\begin{subfigure}[b]{0.34\textwidth}\centering
\begin{tikzpicture}
\begin{axis}[width=5.2cm, height=4.5cm, axis lines=left, xmin=0, xmax=3.5, ymin=0, ymax=1.6,
   xlabel={$r/a$}, xtick={0,1,2,3}, ytick={0,0.5,1,1.5}, grid=major,
   legend style={at={(0.5,1.03)}, anchor=south, legend columns=2, font=\scriptsize}]
  \addplot[curvaE, domain=0:1, samples=2] {x}; \addlegendentry{$E$}
  \addplot[curvaV, domain=0:1, samples=40] {(3-x^2)/2}; \addlegendentry{$\phi$}
  \addplot[curvaE, domain=1:3.5, samples=60, forget plot] {1/x^2};
  \addplot[curvaV, domain=1:3.5, samples=60, forget plot] {1/x};
\end{axis}
\end{tikzpicture}
\caption{$E$ y $\phi$}
\end{subfigure}\hfill
\begin{subfigure}[b]{0.34\textwidth}\centering
\begin{tikzpicture}
\begin{axis}[width=5.2cm, height=4.5cm, axis lines=left, xmin=0, xmax=3.5, ymin=0, ymax=1.15,
   xlabel={$r/a$}, ylabel={$\dd U_e/\dd r$}, xtick={0,1,2,3}, ytick={0,0.5,1}, grid=major]
  \addplot[name path=f1, curvaD, domain=0:1, samples=40] {x^4};
  \addplot[name path=f2, curvaD, domain=1:3.5, samples=60] {1/x^2};
  \addplot[name path=ej, draw=none, domain=0:3.5] {0};
  \addplot[fill=campoD, fill opacity=0.18, draw=none] fill between[of=f1 and ej, soft clip={domain=0:1}];
  \addplot[fill=campoE, fill opacity=0.18, draw=none] fill between[of=f2 and ej, soft clip={domain=1:3.5}];
  \node[etiqt] at (axis cs:0.72,0.12) {$\tfrac16$};
  \node[etiqt] at (axis cs:1.75,0.15) {$\tfrac56$};
\end{axis}
\end{tikzpicture}
\caption{Energía por unidad de radio}
\end{subfigure}
\caption{\Cref{prob:esfera-cargada}. (b) Campo (unidades $Q/4\pi\eo a^2$) y potencial (unidades $Q/4\pi\eo a$): el campo crece linealmente
dentro y decae como $1/r^2$ fuera; el potencial es parabólico dentro, con $\phi(0)=\frac32\phi(a)$. (c) Energía por unidad de radio,
$\dd U_e/\dd r=4\pi r^2u_e$: crece como $r^4$ dentro y decae como $r^{-2}$ fuera. Las áreas sombreadas dan el reparto de la energía:
una sexta parte dentro de la esfera y cinco sextas partes fuera.}
\label{fig:p51}
\end{figure}
